\documentclass[12pt]{article}
\usepackage{amsmath,amssymb,amsfonts}
\usepackage[margin=1in]{geometry}

\begin{document}

\section*{Chapter 7, Problem 6}

\textbf{Problem:} Prove that the Wronskian of a second-order linear differential operator never vanishes. Do this by first proving that $W(x)$ cannot be identically equal to zero (because if it were, $u_1$ and $u_2$ would be linearly dependent, contrary to assumption). Then, differentiating $W(x)$ and using the fact that $u_1$ and $u_2$ are solutions to
\[
f_0(x)\frac{d^2u}{dx^2} + f_1(x)\frac{du}{dx} + f_2(x)u = 0,
\]
prove that
\[
\frac{dW}{dx} = -\frac{f_1}{f_0}\,W.
\]
Integrate this and show that if $f_0(x)$ and $f_1(x)$ are continuous and if $f_0(x)$ never vanishes, then $W(x)$ never vanishes.

\bigskip
\textbf{Solution:}

\textbf{Step 1: $W$ is not identically zero.}

Let $u_1$ and $u_2$ be two linearly independent solutions of the homogeneous equation. The Wronskian is:
\[
W(x) = u_1(x)u_2'(x) - u_2(x)u_1'(x).
\]

If $W(x) = 0$ for all $x$, then $u_1 u_2' = u_2 u_1'$. At any point where $u_1 \neq 0$:
\[
\frac{d}{dx}\Bigl(\frac{u_2}{u_1}\Bigr) = \frac{u_1 u_2' - u_2 u_1'}{u_1^2} = \frac{W}{u_1^2} = 0,
\]
so $u_2/u_1 = \text{const}$, meaning $u_2 = c\,u_1$. This contradicts the assumption that $u_1$ and $u_2$ are linearly independent. Hence $W$ is not identically zero.

\textbf{Step 2: Derive the ODE for $W$.}

Since $u_1$ and $u_2$ both satisfy the equation:
\[
f_0\,u_i'' + f_1\,u_i' + f_2\,u_i = 0 \quad \Longrightarrow \quad u_i'' = -\frac{f_1}{f_0}\,u_i' - \frac{f_2}{f_0}\,u_i, \quad i = 1,2.
\]

Differentiate the Wronskian:
\begin{align}
W' &= u_1' u_2' + u_1 u_2'' - u_2' u_1' - u_2 u_1'' \\
&= u_1 u_2'' - u_2 u_1''.
\end{align}

Substitute the expressions for $u_1''$ and $u_2''$:
\begin{align}
W' &= u_1\Bigl(-\frac{f_1}{f_0}u_2' - \frac{f_2}{f_0}u_2\Bigr) - u_2\Bigl(-\frac{f_1}{f_0}u_1' - \frac{f_2}{f_0}u_1\Bigr) \\
&= -\frac{f_1}{f_0}(u_1 u_2' - u_2 u_1') - \frac{f_2}{f_0}(u_1 u_2 - u_2 u_1) \\
&= -\frac{f_1}{f_0}\,W.
\end{align}

\textbf{Step 3: Integrate and conclude.}

This is Abel's identity. The ODE $W' = -(f_1/f_0)\,W$ is separable:
\[
\frac{dW}{W} = -\frac{f_1(x)}{f_0(x)}\,dx.
\]

Integrating from $x_0$ to $x$:
\[
W(x) = W(x_0)\exp\Bigl(-\int_{x_0}^x \frac{f_1(t)}{f_0(t)}\,dt\Bigr).
\]

Since we showed $W$ is not identically zero, there exists a point $x_0$ where $W(x_0) \neq 0$. The exponential function is never zero, and the integral $\int_{x_0}^x f_1/f_0\,dt$ is finite as long as $f_0$ and $f_1$ are continuous and $f_0(x) \neq 0$ on the interval.

Therefore $W(x) \neq 0$ for all $x$ in the interval. $\blacksquare$

\end{document}
